% Created 2026-06-06 Sat 17:22
% Intended LaTeX compiler: pdflatex
\documentclass[11pt]{article}
\usepackage[utf8]{inputenc}
\usepackage[T1]{fontenc}
\usepackage{graphicx}
\usepackage{longtable}
\usepackage{wrapfig}
\usepackage{rotating}
\usepackage[normalem]{ulem}
\usepackage{amsmath}
\usepackage{amssymb}
\usepackage{capt-of}
\usepackage{hyperref}
\date{2026-06-06}
\title{Preset UX And Wick Lowering Improvement Plan}
\hypersetup{
 pdfauthor={},
 pdftitle={Preset UX And Wick Lowering Improvement Plan},
 pdfkeywords={},
 pdfsubject={},
 pdfcreator={Emacs 30.2 (Org mode 9.7.11)}, 
 pdflang={English}}
\begin{document}

\maketitle
\tableofcontents

\section{Goal}
\label{sec:org71835f3}

This note joins the deficiencies of the current \href{../org/cft-code/presets/presets.org}{Presets} skeleton into one
coherent improvement plan.  The current code gives us a checked public shape,
but it is not yet the right user experience and it does not encode real Wick
rules.

The target is:

\begin{itemize}
\item users write physics objects, not ids;
\item operator builders preserve labels;
\item presets lower to \href{../org/cft-code/theory/theory.org}{Theory} data;
\item \href{../org/cft-code/correlators/wick.org}{Wick} rules and zero-mode rules are compact generated tables;
\item \href{../org/cft-code/correlators/correlators.org}{Correlators} still receive only local operators, a compiled config, and a
sink.
\end{itemize}
\section{Current Deficiencies}
\label{sec:orge1366ab}

The present implementation has six important gaps.

\begin{enumerate}
\item \texttt{Handle.Index}, \texttt{Handle.Momentum}, \texttt{Handle.Profile}, and related values are
naked integer aliases.  That leaks storage details into the physics API.
\item Operator builders discard the labels they receive.  \texttt{expX(k,z,zb)} and
\texttt{expX(k2,z,zb)} currently produce the same body id.
\item \texttt{config.sphere} and \texttt{config.disk} contain placeholder Wick and zero-mode ids.
\item All preset \texttt{correlator} functions return \texttt{error.NotImplemented}.
\item \texttt{product} and \texttt{boundary} expose hard-coded wrapper builders for the current
milestone rather than selecting namespaces from supplied factors.
\item Boundary extensions are accepted as data, but the generated boundary builder
does not yet enforce which extensions are present.
\end{enumerate}

The core mistake is not that ids exist.  The mistake is that the user-facing
surface suggests users might type those ids by hand.  Ids should be assigned by
generated stores and used only by lowering/runtime code.
\section{UX Boundary}
\label{sec:orgb013518}

There should be three layers.

\begin{enumerate}
\item The user-facing physics layer has typed values such as \texttt{Index}, \texttt{Momentum},
\texttt{Profile}, \texttt{TargetPoint}, and \texttt{BoundaryStack}.  These are opaque structs, not
aliases.
\item The generated builder layer interns those values into a compact label store
and builds local operators.
\item The lowered runtime layer sees only small ids, label spans, Wick-rule tables,
zero-mode tables, and coefficient streams.
\end{enumerate}

The user should write code like this:

\begin{verbatim}
const X = cft.presets.freeBoson(.{ .dimension = 10, .alpha_prime = .one });

var local = X.local(allocator);
const mu = try local.index("mu");
const k = try local.momentum("k");
const z = try local.coord("z");
const zb = try local.coord("zb");

const ops = try local.ops(.{
    local.op.dX(mu, 0, z),
    local.op.expX(k, z, zb),
});

try X.correlator(X.config.sphere, ops, sink);
\end{verbatim}

This does not introduce a correlator context.  \texttt{local} is a short-lived builder
and label arena for the operator collection.  The correlator still receives a
config, a collection of local operators, and a sink.
\section{Opaque Physics Values}
\label{sec:orgc5ce9e7}

The public aliases in \texttt{Handle} should be replaced by generated typed tokens.
The internal integer remains, but ordinary user code never needs to type it or
read it.  In Zig this can be represented as a non-exhaustive integer-backed
enum, with private generated helpers for extraction.

\begin{verbatim}
pub const Index = enum(u32) { _ };
pub const Momentum = enum(u32) { _ };
pub const Profile = enum(u32) { _ };
pub const TargetPoint = enum(u32) { _ };
pub const BoundaryStack = enum(u32) { _ };
\end{verbatim}

These tokens are cheap to copy and easy for the optimizer.  They also prevent
passing a momentum where an index is required.

For boundary data, the user-facing shape should be similarly physical:

\begin{verbatim}
const brane = cft.presets.freeBosonBoundary(.{
    .neumann = X.target.subspace(.{ 0, 1, 2, 3 }),
    .dirichlet = X.target.complement(.{ 0, 1, 2, 3 }),
    .dirichlet_position = try local.point("x0"),
    .chan_paton = cp.stack("N"),
});
\end{verbatim}

The projectors and Chan-Paton tensors may lower to \href{../org/tensor-code/tensor-code.org}{Tensor-Code} handles, but
the user should not manually provide those handles in ordinary code.
\section{Local Builder And Label Store}
\label{sec:orga9187b6}

The generated \texttt{local} builder owns a compact label arena.  Operator builders
write label records to that arena and return local operators that point to
label spans.

\begin{verbatim}
const LocalOp = struct {
    insertion: OperatorInsertion,
    kind: OperatorKindId,
    labels: LabelSpan,
};

const MultiOp = struct {
    operators: []const LocalOp,
    labels: *const LabelStore,
};
\end{verbatim}

This keeps the runtime call narrow while preserving labels.  A \texttt{MultiOp} is
still just a collection of local operators; it carries the label store because
the collection is not meaningful without the labels.

The builder should append small fixed-width label records.  It should not
allocate expression trees.

\begin{verbatim}
const LabelValue = union(enum) {
    integer: i64,
    rational: Rational,
    tensor: tensor.kernel.TensorExprId,
    symbol: SymbolId,
};
\end{verbatim}

For common labels such as vector indices, derivative labels, momenta, and
profile ids, the records are tiny.  Larger physics objects remain behind
handles owned by the preset, tensor-code, or the client.
\section{Operator Builders}
\label{sec:org5e49211}

Operator builders should stop returning a kind-only body.  They should emit a
local operator with the appropriate label span.

\begin{verbatim}
pub fn expX(self: *Local, k: Momentum, z: Coord, zbar: Coord) !LocalOp {
    const labels = try self.labels.append(.{ .momentum = raw(k) });
    return self.op(.exp_x, .bulk(z, zbar), labels);
}
\end{verbatim}

The same pattern covers \texttt{dX(mu,n,z)}, \texttt{profile(f,z,zb)},
\texttt{dXBoundary(mu,n,y)}, and \texttt{cBoundary(n,y)}.  The generated builder knows the
schema; the user only sees the typed arguments.
\section{Wick Rule Encoding}
\label{sec:org9c120bc}

The current code does not encode Wick rules.  The replacement should be a
sparse rule table per \texttt{CorrelatorConfig}.

\begin{verbatim}
const WickRule = struct {
    left: OperatorKindId,
    right: OperatorKindId,
    evaluator: WickEvaluator,
    data: RuleDataId,
};

const WickRuleSet = struct {
    rules: []const WickRule,
};
\end{verbatim}

The evaluator should be either a small generated tag for built-in fast paths or
a client-supplied procedure for genuinely custom theories.

\begin{verbatim}
const WickEvaluator = union(enum) {
    generated: GeneratedWickEvaluator,
    custom: CustomWickEvaluator,
};
\end{verbatim}

The evaluator receives only what it needs: two local operators, the label
store, rule data, and an output sink.

\begin{verbatim}
fn evalPair(
    rule_data: RuleDataId,
    left: LocalOp,
    right: LocalOp,
    labels: *const LabelStore,
    sink: anytype,
) !void;
\end{verbatim}

For the first presets, the generated evaluator tags should cover:

\begin{itemize}
\item free-boson derivative/derivative pairings;
\item free-boson derivative/exponential pairings;
\item free-boson exponential/exponential multiplicative factors;
\item profile rewrites before Wick evaluation;
\item \(bc\) first-order pairings;
\item boundary free-boson and mixed bulk-boundary kernels;
\item boundary \(bc\) and mixed bulk-boundary \(bc\) kernels.
\end{itemize}
\section{Zero-Mode Rules}
\label{sec:org3fcd188}

Zero-mode rules should be separate from pairwise Wick rules.  They consume the
residual uncontracted local operators after Wick recursion.

\begin{verbatim}
const ZeroModeRule = struct {
    sector: SectorId,
    evaluator: ZeroModeEvaluator,
    data: RuleDataId,
};

const ZeroModeRuleSet = struct {
    rules: []const ZeroModeRule,
};
\end{verbatim}

For the first presets, the generated zero-mode evaluators should cover:

\begin{itemize}
\item free-boson noncompact momentum conservation;
\item explicit zero-mode integrals for position-space profiles;
\item Neumann projected momentum conservation;
\item Dirichlet evaluation at the brane position;
\item sphere \(bc\) top-form saturation;
\item disk boundary \(bc\) top-form saturation.
\end{itemize}

Momentum deltas, profile integrals, and brane-position phases should be
coefficient factor families.  They should not add new branches to the
correlator algorithm.
\section{Correlator Config}
\label{sec:orgd717ccc}

A config should point to actual tables, not placeholder ids.

\begin{verbatim}
pub const CorrelatorConfig = struct {
    wick_rules: WickRuleSetId,
    zero_modes: ZeroModeRuleSetId,
};
\end{verbatim}

The runtime call remains:

\begin{verbatim}
try Theory.correlator(Theory.config.sphere, ops, sink);
\end{verbatim}

Here \texttt{ops} already carries its label store.  The config chooses the Wick and
zero-mode tables.  No surface object or boundary-condition object is passed at
runtime.
\section{Product And Boundary Composition}
\label{sec:orgd61e59b}

The current hard-coded product/boundary wrapper is acceptable as a temporary
scaffold, but the improved lowering should merge generated namespaces from the
supplied factors and extensions.

The rule is simple:

\begin{itemize}
\item \texttt{product(.\{ X, bc \})} exposes only builders contributed by \texttt{X} and \texttt{bc};
\item \texttt{boundary(Bulk, .\{ freeBosonBoundary(...) \})} exposes only the selected
boundary free-boson builders;
\item adding \texttt{bcDiskBoundary(.\{\})} exposes boundary ghosts and mixed ghost rules;
\item Chan-Paton stacks remain labels on boundary operators and may point to
tensor-code data.
\end{itemize}

This avoids a generic runtime registry.  Composition is a compile-time lowering
step over a small tuple of factors.
\section{Implementation Order}
\label{sec:orgc2c2dcb}

\begin{enumerate}
\item Replace public naked \texttt{Handle} aliases in \href{../org/cft-code/presets/presets.org}{Presets} with generated opaque
physics structs.
\item Add a generated \texttt{Local} builder with a compact label arena and typed
constructors for coordinates, indices, momenta, profiles, target points,
subspaces, and Chan-Paton stacks.
\item Change operator builders to preserve label spans in \texttt{LocalOp}.
\item Replace placeholder \texttt{CorrelatorConfig} ids with real \texttt{WickRuleSet} and
\texttt{ZeroModeRuleSet} ids.
\item Implement generated Wick evaluator tags for free boson and \(bc\) on the
sphere.
\item Implement generated zero-mode evaluators for free-boson momentum/profile
rules and \(bc\) top forms.
\item Extend the same lowering to disk free-boson BCFT and boundary \(bc\).
\item Replace hard-coded product/boundary wrappers with compile-time namespace
composition over supplied factors.
\end{enumerate}
\section{Tests Worth Keeping}
\label{sec:org017d4c1}

Tests should check design invariants, not syntax.

\begin{itemize}
\item \texttt{expX(k1,...)} and \texttt{expX(k2,...)} preserve distinct momentum labels.
\item \texttt{dX(mu,n,z)} preserves both the vector index and derivative label.
\item \texttt{product(.\{X,bc\})} exposes exactly the expected bulk builders.
\item \texttt{boundary(...freeBosonBoundary...)} exposes boundary free-boson builders and
does not expose boundary ghosts unless \texttt{bcDiskBoundary} is present.
\item sphere \(bc\) top form evaluates the three-\(c\) correlator normalization.
\item free-boson sphere zero modes emit a momentum-conservation factor for plane
waves.
\end{itemize}

The first two tests should be written before any Wick evaluator.  They protect
the UX/data boundary that is currently broken.
\section{What Not To Do}
\label{sec:orgd4175c1}

Do not add a general runtime "context" object to correlator calls.  The operator
collection can carry labels; the config can carry Wick and zero-mode tables.
That is enough.

Do not make coefficient-code learn about every physics case.  Momentum deltas,
profile integrals, Chan-Paton traces, pure-spinor measures, and lattice theta
factors should remain factor families.

Do not implement a broad plugin registry or reflection-heavy namespace merger
before the first free-boson/\(bc\) presets work.  The compile-time tuple of
factors is sufficient for now.

Do not make Wick rules allocate full expanded expressions.  Rule evaluators
stream coefficient factors and residual local operators into sinks.
\section{AGENTS Audit}
\label{sec:orgc7e7f0d}

This plan is aligned with \texttt{AGENTS.md}.

\begin{itemize}
\item Data-oriented: labels are compact arena records; local operators point to
spans; Wick and zero-mode rules are sparse tables.
\item Minimal runtime data: correlator receives a config, a local-operator
collection, and a sink.
\item Memory discipline: rule evaluators stream terms and do not materialize large
symbolic trees.
\item Parallelizable flow: sectors and independent Wick branches can be streamed
independently once the label store is immutable.
\item Minimal public API: users see typed physics builders; ids, label spans,
evaluator ids, and rule data stay private.
\item Literate structure: this note links to the relevant literate modules; the
implementation should happen in \href{../org/cft-code/presets/presets.org}{Presets}, \href{../org/cft-code/theory/theory.org}{Theory},
\href{../org/cft-code/kernel/cft-kernel.org}{CFT Kernel}, \href{../org/cft-code/correlators/wick.org}{Wick}, and
\href{../org/cft-code/correlators/correlators.org}{Correlators}.
\item Test discipline: the proposed tests protect label preservation, namespace
selection, zero-mode rules, and first real physics checks.
\end{itemize}

The one design point that needs explicit sign-off is whether \texttt{MultiOp} carrying
a pointer to its label store is acceptable terminology.  It preserves the
principle that correlator does not receive a separate context, but it makes the
operator collection self-contained enough for fast runtime evaluation.
\end{document}
